\documentclass[runningheads]{llncs}
\usepackage[T1]{fontenc}
\usepackage{graphicx}
\usepackage[english]{babel}
\usepackage{makecell}
\usepackage{listings}
\usepackage{booktabs}
\usepackage{tabularx}
\usepackage[disable]{todonotes}

\renewcommand{\cellalign}{l}

\begin{document}
\title{Lightweight confidential and traceable multi-stakeholder data processing}

\author{Steylor Barros\inst{1}\orcidID{?????} \and
Eduardo Falcão\inst{2}\orcidID{0000-zzz-zzz-zzzz} \and Andrey Brito\inst{1}\orcidID{0000-0002-7350-8599}}
\authorrunning{S. Barros et al.}

\institute{Universidade Federal de Campina Grande, Campina Grande/PB, Brazil \and Universidade Federal do Rio Grande do Norte, Natal/RN, Brazil 
\email{steylor.barros@ccc.ufcg.edu.br}\\
\email{eduardo@dca.ufrn.br}\\
\email{andrey@computacao.ufcg.edu.br}\\
\url{https://www.lsd.ufcg.edu.br/}}

\maketitle

\begin{abstract}
% [AB] Essa é uma primeira versão para organizar meus pensamentos. Ainda vai ser melhorada.
We leverage technologies such as confidential computing and distributed ledgers to propose an architecture for multi-stakeholder, confidential data processing. Our approach aims to protect non-interactive data-processing applications in scenarios where multiple parties (e.g., data controllers, application developers, and infrastructure managers) distrust one another. We generalize previous results, describe the architecture and requirements of the involved components, and present an implementation focused on lightweight execution in the AWS cloud. Our implementation focuses on Python applications, showing overheads as low as $X\%$.

% TODO: [AB] seria bom ter uma applicação realista. 

\keywords{Device Attestation, User Authentication, Zero Trust Security,
Virtual Private Network (VPN), Remote Attestation, SPIFFE, SPIRE,
Multi-Factor Authentication (MFA)}
\end{abstract}

\section{Introduction}

% [AB]: todo
% Falar mais do aspecto de multi-stakeholder
% Ele cita o DNAT mas deve comparar melhor, mostrando que tem limitações (ex., a execução em SGX dificulta o reuso de aplicações, que podem precisar serem recompiladas ou adaptadas, e também tem o problema das dependências).

Confidential data marketplaces must simultaneously support asset
exchange, access control, and third-party computation over sensitive
inputs. That combination creates a difficult systems problem.
A marketplace that accepts externally supplied applications cannot
treat code execution as an isolated implementation detail, because
the same platform also handles persistent metadata, payment and
access logic, storage backends, and the datasets that motivate the
marketplace in the first place. In practice, the problem is not merely
how to ``run code safely,'' but how to prevent build-time and
run-time compromise from expanding into a full platform compromise.

The most direct implementation path would be a centralized
confidential virtual machine~\cite{amd2020sevsnp} that hosts the
marketplace-facing logic, local artifact storage, and the execution
orchestrator in the same trust domain. That design is operationally
attractive because it minimizes deployment complexity and avoids
inter-service coordination. However, it also creates the wrong
security shape for the problem: the same node that validates access
and manages persistent sensitive state also becomes the place where
untrusted or semi-trusted workloads are prepared and executed.

This paper argues that such a concentration of privilege is
fundamentally ill-suited for high-assurance marketplaces. It is
proposed an architecture designed under the assumption that component
compromise is an eventual certainty rather than a remote possibility.
By decomposing the marketplace into distinct trust
domains---control, build, and execution~\cite{nascimento2020dnat}---the
system ensures that a failure or malicious exploit in one phase cannot
propagate to the entire application. Our focus is on minimizing the
``blast radius'' of errors and attacks, ensuring that even if the
build or execution planes are corrupted, the platform's core sensitive
assets and persistent state remain protected.


% qual é o contexto do problema?

% por que não é trivial?

% O que existe por aí?
%% menção rápida ao dnat e suas limitações 

% Uma possível solução

% Limitações 
% - da arquitetura (ex. execução offline)
% - da implementação (ex., uso de python)

% organização do artigo

\section{Background}\label{sec:background}

% [AB] Aqui em background, deveríamos:
% (1) explicar melhor o que é oportunidade trazida pela computação confidencial: não precisa copiar na infraestrutura e simplifica o papel de validação da aplicação. 
% (2) explicar blockchain e como ela permite o rastreamento;
% (3) explicar por que IPFS resolve limitações da blockchain.
% Importante: agora que Steylor está avançado e já sabe o que vai implementar e o que não vai implementar, podemos evitar falar das coisas que ele não vai de fato usar. Por exemplo, será que precisa de IPFS? Não é melhor falar que precisa de um armazenamento distribuído compatível com os objetivos e depois dizer que IPFS é uma opção e que aqui estamos unsando ele de forma simplificada?

\subsection{Confidential computing}

%[AB] Acho que essa seção está abordando o ponto errado. Nossa ênfase não é computação confidencial, então podemos dizer o que de computação confidencial

Confidential Computing~\cite{amd2020sevsnp} is a paradigm that focuses on protecting
data while it is being processed. It leverages hardware-level memory encryption to make
a process or an entire Virtual Machine inaccessible to the host hypervisor. By using an
ASID (Address Space Identifier), the processor applies cryptography to a memory context;
the ciphertext is decrypted transparently by a key stored directly in the memory
controller when the VM that owns that context accesses the data. Confidential Virtual
Machines are recognized by a special bit, called the C-bit, which instructs the CPU to
restrict access to the corresponding memory pages. Any attempt by the hypervisor to map
those pages into its own address space raises a processor exception, enforcing
hardware-backed isolation even against a privileged adversary.

\subsection{Blockchain}
A blockchain~\cite{nakamoto2008bitcoin} is a peer-to-peer distributed system that
records transactions with strong integrity guarantees. Transactions are grouped into
blocks; each block contains a cryptographic hash of the previous block, creating a chain
in which every entry implicitly attests to the entire preceding history. Altering any
block invalidates all subsequent hashes, so corrupting the ledger requires recomputing
more than fifty percent of the combined processing capacity of all participating
peers---an effort generally considered computationally infeasible in practice.
Blockchains also support Smart Contracts~\cite{buterin2014ethereum}, which are programs
deployed and executed directly on the network. A Smart Contract can register and govern
access to relevant data, such as the encryption keys and access permissions used in our
system. The main drawback of blockchains is their inherent transaction cost: every write
operation requires consensus among the majority of peers before a block is accepted,
which introduces both financial cost and latency.

\subsection{IPFS}
The InterPlanetary File System (IPFS)~\cite{benet2014ipfs} addresses the challenge of
storing large data volumes that do not integrate well with a blockchain due to the costs
described above. IPFS is a structured distributed system that uses a Distributed Hash
Table (DHT) to locate files across the network. Files are identified by their content
hash---a deterministic, one-way function that produces a unique identifier for a given
sequence of bytes; any modification to the content produces a different hash. The DHT
ensures that data remains accessible to all peers by managing replication and routing
efficiently across the network, making IPFS a practical complement to blockchain-based
access control for large binary artifacts.



% -------------------------------------------------------------
\section{Problem definition and solution requirements}
\label{sec:problem-definition}

The core challenge addressed by the Dataset and Application Marketplace (DNAT) involves
a transaction between two mutually distrusting parties: a \textit{data buyer}, who owns
a confidential application, and a \textit{dataset provider}, who holds sensitive data.
The objective is to enable a computation in which the buyer's application processes the
provider's dataset without exposing either asset to the other party or to the underlying
cloud infrastructure.

A straightforward implementation would deploy a single Confidential Virtual Machine (CVM)
responsible for all platform functions: receiving uploads, managing the IPFS storage layer,
interacting with the Smart Contract, resolving and installing Python dependencies, and
running the application against the dataset. While operationally simple, this monolithic
design concentrates all privileges within one trust domain. The critical vulnerability
emerges during dependency resolution: installing external packages to satisfy an
application's requirements necessarily executes arbitrary third-party code with the same
privileges as the marketplace layer, the IPFS node, and the blockchain client. A
maliciously crafted package---or a compromised dependency in the supply chain---could
exploit this phase to exfiltrate stored datasets, tamper with marketplace metadata, or
establish a persistent backdoor affecting future executions. A build-time compromise thus
propagates immediately into a full platform compromise, making dependency resolution a
disproportionately dangerous entry point.

This observation reframes the core problem. Secure execution of the application at runtime
is necessary but not sufficient. An equally important concern is preventing a compromise
during the build and dependency resolution phase from reaching the broader platform.
The problem must therefore be treated as a \textit{privilege separation problem}: distinct
operational phases---control, build, and execution---carry fundamentally different risk
profiles and must occupy separate trust domains with different network-access policies,
persistence rights, and failure blast radii.

Three design consequences follow from this reframing. First, the component managing
marketplace state, blockchain interaction, and IPFS operations must never directly execute
dependency installation logic for user-supplied applications. Second, the build phase,
where external network access is unavoidable, must be isolated from both persistent
sensitive storage and the execution runtime. Third, the execution phase, during which the
sensitive dataset is consumed, must operate without any outbound network interface,
eliminating the primary exfiltration channel even under a complete runtime compromise.

A secure solution must satisfy the following requirements:

% Multi-stakeholder machine learning

% Termina com uma lista de requisitos ou desafios.

\section{Proposal}\label{sec:proposal}

\subsection{Architecture}
The DNAT platform is organized around three Confidential Virtual
Machines, each assigned a distinct operational role and a distinct
set of persistence rights. \textbf{CVM1} acts as the control plane,
\textbf{CVM2} as the isolated build plane, and \textbf{CVM3} as the
isolated execution plane. The central design principle is privilege
separation: a compromise in one plane must not grant access to the
assets of the others. (...)

% conecta tecnologias-chave para resolver o problema

% como cada tecnologia ajuda a resolver os desafios/requisitos

\subsection{Implementation}
The architecture described above is realized using Firecracker
microVMs~\cite{agache2020firecracker} as the primary isolation
primitive for both the build and execution environments. The current
prototype instantiates each CVM as a containerized service for
development purposes; in a production deployment, each service would
run inside a hardware-backed Confidential VM~\cite{amd2020sevsnp}. (...)

% =====================================================================
% Secao de Implementation (nova). Descreve o sistema concreto, fundamentada
% no codigo do repositorio. Idioma: ingles. Sugestao: inserir logo apos a
% Secao "DNAT Architecture" e antes da "Security Analysis".
% =====================================================================

\section{Implementation}
\label{sec:implementation}

This section describes how the three planes are realized, so that the security properties
evaluated later can be traced to concrete mechanisms rather than to design intent alone. The
system is deployed as three containers on a shared runtime---\texttt{dnat-client} (CVM1),
\texttt{dnat-builder} (CVM2), and \texttt{dnat-executor} (CVM3)---each standing in for the
hardware-backed confidential VM~\cite{amd2020sevsnp} of a production deployment, and is
contrasted with two baselines: a single-container deployment that colocates the same assets
(A) and a no-isolation variant that removes the microVM entirely (M0). The consequences of
that substrate substitution are examined in Section~\ref{sec:evaluation}.

\subsection{Control Plane (CVM1)}

The control plane hosts the marketplace API, a local IPFS node, and the blockchain client,
exposing registration, access checks, and execution orchestration through a small HTTP API.
On registration, CVM1 does not install dependencies itself: it assembles a build bundle (the
source plus a manifest of declared dependencies) and forwards it to the builder. When the
builder returns an \texttt{application.ext4} artifact, CVM1 encrypts it with AES-256-GCM
over a gzip-compressed image, in a versioned envelope carrying the IV, authentication tag,
and metadata. Only the encrypted blob and a manifest are published to IPFS; the plaintext
artifact never leaves the control plane. Access control is enforced on chain: an execution
request first calls \texttt{hasAccessByIds}, and the workflow aborts with \emph{access
denied} if the caller has not purchased access to both assets. A purchased right is recorded
under the key \texttt{keccak256(encryptedDatasetHash, encryptedApplicationHash,
buyer)}---bound to the encrypted IPFS CIDs of the two assets and the buyer's address rather
than to mutable identifiers, so the grant is stable under off-chain re-indexing---and is
verified before any bundle is assembled.

\subsection{Build Plane (CVM2)}

The builder exposes a single \texttt{/build} endpoint that accepts a bundle and launches an
ephemeral Firecracker microVM to resolve dependencies. Unlike the executor, the build
microVM is given a network interface, because resolving and compiling Python wheels may
require egress to package indexes. Egress is constrained on the host side: a TAP device is
created and iptables rules \texttt{REJECT} traffic from the guest toward private ranges
(\texttt{10.0.0.0/8}, \texttt{172.16.0.0/12}, \texttt{192.168.0.0/16},
\texttt{127.0.0.0/8}, \texttt{169.254.0.0/16}, and the host TAP address), while the
remaining outbound traffic is masqueraded. Inside the guest, an artifact builder produces a
read-only \texttt{ext4} image holding the application source, its \texttt{site-packages}, a
lock file of resolved dependencies, and a marker file used later for disk discovery.
Dependencies are first wheel-built and then installed with \texttt{--no-index} from the
local wheelhouse; reusable wheels are persisted in a size-pruned cache that is the builder's
only durable state---and, as Section~\ref{sec:evaluation} notes, its only cross-request
channel. After the build, the host returns only the artifact and its metadata; the
per-request build state is discarded.

\subsection{Execution Plane (CVM3)}

The executor exposes a single \texttt{/execute} endpoint. For each request it copies a fresh
root overlay, creates new input and output \texttt{ext4} disks, optionally attaches the
application artifact read-only, and launches a Firecracker microVM with one vCPU, 512~MiB of
memory and---crucially---no network interface. The guest does not receive disk roles out of
band; it discovers them by scanning candidate block devices for marker files, then extracts
the workspace and runs \texttt{workspace/run.sh}, whose output is persisted to the output
disk. Host and guest synchronize over the serial console: the guest writes an
\texttt{EXECUTION\_COMPLETE} marker, after which the host issues a graceful shutdown,
escalating to termination if needed, then mounts the output disk and derives the HTTP
response solely from the persisted result. A \texttt{cleanup} trap unmounts every loop device
and removes the per-run working directory on exit, which is what makes execution state
ephemeral.

\subsection{Application Contract, Asset Integrity, and Baselines}

An application is registered as unmodified source together with a manifest of its
third-party dependencies, which the builder resolves and bundles into the read-only
\texttt{application.ext4} artifact. Because dependencies are shipped as ordinary
\texttt{site-packages} inside that artifact---rather than linked into a trusted-execution
enclave at build time---an application that follows this contract runs with arbitrary
declared dependencies, unchanged across all three environments and without recompilation or
enclave-specific adaptation (R6). This is the reuse property that
Section~\ref{sec:evaluation} contrasts with the enclave-based original DNAT.

Assets are bound to the chain by content hash: registration records the SHA-256 of the
artifact, which lets a consumer detect tampering between IPFS and the contract. The first
baseline (environment~A) runs the identical control, build, and execution logic, including
Firecracker isolation for both microVMs, but inside a single container whose filesystem
holds the IPFS repository, the wheel cache, and the executions store, and whose environment
holds the asset-encryption key and the wallet private key. Its only structural difference
from the compartmentalized deployment is the colocation of these assets---exactly the
variable the security evaluation isolates.

The second baseline removes isolation altogether. With \texttt{DNAT\_NO\_MICROVM} set, the
executor runs the same \texttt{workspace/run.sh} contract directly in the container's
namespace---no Firecracker, no per-run disks, no network removal---while emitting the same
result, so the request path is otherwise unchanged. In this variant (M0) the application
inherits the container's network, environment (including the control-plane secrets), and
persistent filesystem, which is what lets the evaluation separate the microVM's operational
cost from the containment it provides.

% fala de abordagem do build (forma de incorporar as dependências), de microVM e a execução (isolamento de rede, transferência de resultados),

% ------------------------------------------------
% ------------------------------------------------
\section{Evaluation}
\label{sec:evaluation}

% [AB] sugeriu STRIDE para organizar a avaliacao de seguranca. Adotado: a
% subsecao "Security Evaluation with STRIDE" em section_evaluation.tex mapeia
% cada letra do STRIDE para um mecanismo do design e o teste que o evidencia.
% A matriz comparativa entre arquiteturas (M0/A/B/SGX) esta em results_tables.tex
% (Table~\ref{tab:comparison}); celulas TBD sao medidas/citadas no proximo ciclo.

% =====================================================================
% Evaluation: methodology + STRIDE security + operational cost + threats.
% Sprint 1 (consolidation): security reorganized with STRIDE (per [AB]); the
% four-architecture comparison (M0/A/B/SGX) is introduced and summarized in the
% comparative matrix (Table~\ref{tab:comparison}). Cells marked TBD there are
% measured/cited in the follow-up campaign (M0, boot breakdown, AWS cost, gas).
% Language: English. Tables: results_tables.tex + section_results_ab.tex.
% =====================================================================

The architectural argument of this paper is comparative: the compartmentalized
design should contain compromise better than less-separated alternatives while
paying an acceptable operational cost. We evaluate the implementation along two
axes. For \emph{security}, we organize the analysis with the STRIDE
model~\cite{shostack2014threat}, applied per interaction over the elements of
Figure~\ref{fig:architecture}, and show how the design neutralizes---or
explicitly fails to neutralize---each non-trivial threat, mapping every claim
to a concrete test. For \emph{cost}, we measure the latency, build
time, and resource price of the isolation strategy.

The comparison spans up to four points on an isolation spectrum, summarized in
Table~\ref{tab:comparison}: a no-isolation monolith (M0), a centralized
deployment that adds ephemeral microVMs but colocates all assets (Environment~A),
the compartmentalized design that also separates the planes (Environment~B), and,
as a cited reference, the SGX-based original DNAT~\cite{nascimento2020dnat}.
Because A and B share the same Firecracker microVMs, any cost or security
difference between them isolates the effect of \emph{compartmentalization}; M0
isolates the cost and the security delta of the \emph{microVM} primitive itself;
and the SGX reference frames the reuse and dependency limitations that motivate
our change of technology.

\subsection{Methodology}

We exercise the running prototype with a small suite of purpose-built
applications, each designed to attack one STRIDE category, over a synthetic,
PII-free dataset. Tests that target the execution plane submit an application
bundle directly to the executor, reproducing the same execution contract the
control plane generates, so that the variable under study---the execution
microVM---is isolated from unrelated services; tests that target lateral movement
inspect the executor environment, and the access-control test drives the
marketplace API. Each adverse application emits machine-readable evidence markers
that a harness collects, and each test is mapped to the interaction and STRIDE
category it supports (Table~\ref{tab:stride}), so that every observation
connects to a property of a specific boundary crossing rather than merely
listing successful checks.

Environments~A and~B were both measured live and end to end. The centralized
baseline (Environment~A) keeps execution under the same Firecracker isolation as
the compartmentalized design (Environment~B), so the only independent variable
between the two is the separation of planes; any measured difference therefore
reflects compartmentalization, not a weaker execution sandbox. The no-isolation
baseline (M0)---which runs the application directly in the control namespace,
without launching a microVM---was measured on the same host, and the on-chain gas
cost was measured on Hardhat; together these complete the comparative matrix
(Table~\ref{tab:comparison}), leaving only the SGX reference column, which is
cited from the original DNAT work.

\subsection{Security Evaluation with STRIDE}

STRIDE is an analysis of components and their interactions, not of a system
taken as a whole: a claim such as ``the design resists tampering'' is only
meaningful relative to a stated element. We therefore fix the level of
abstraction first. The analysis is performed over the elements of
Figure~\ref{fig:architecture}---the external parties, the three planes (each a
separate trust domain), the ledger, the store, and the package indexes---and
over the numbered interactions I1--I8 that cross their trust boundaries.
Following STRIDE-per-element practice~\cite{shostack2014threat}, each threat
category was examined against every element and boundary-crossing interaction;
flows internal to a single trust domain are not analyzed.
Table~\ref{tab:stride-matrix} records the full sweep---every category against
every interaction---classifying each cell as actively countered, trivially
handled, or an acknowledged residual, so that the coverage can be audited
rather than taken on trust. Table~\ref{tab:stride} then concentrates on the
non-trivial threats---those the architecture must actively counter---naming for
each the exact interaction where it arises, the countering mechanism, and the
test that provides evidence. The trivial and residual cells are discussed in
the text at the end of this subsection.

\begin{table}[htbp]
\centering
\caption{Completeness matrix: every STRIDE category against every interaction
of Figure~\ref{fig:architecture} (last two rows are element-level attack
surfaces, not legitimate flows). $\bullet$~countered by a design mechanism
(detailed in Table~\ref{tab:stride} or the text); $\circ$~trivial---standard
TLS/HTTPS or a trust assumption of Section~\ref{sec:proposal};
$\odot$~acknowledged residual or out of scope, discussed in the text and
Threats to Validity; ``--''~not applicable to that element or flow, or the
availability of external infrastructure (ledger, package indexes) assumed out
of scope. MicroVM/VMM escape at I7 is excluded by the trusted-VMM assumption of
Section~\ref{sec:proposal}; the reverse-pivot row captures what a plane
compromise reaches despite it.}
\label{tab:stride-matrix}
\footnotesize
\begin{tabular}{@{}lcccccc@{}}
\toprule
Interaction / element & S & T & R & I & D & E \\
\midrule
I1\quad external party $\leftrightarrow$ control & $\bullet$ & $\circ$ & $\bullet$ & $\circ$ & $\odot$ & $\circ$ \\
I2\quad control $\leftrightarrow$ ledger & $\circ$ & $\circ$ & $\bullet$ & $\circ$ & -- & -- \\
I3\quad control $\leftrightarrow$ store & $\circ$ & $\bullet$ & -- & $\bullet$ & -- & -- \\
I4\quad control $\rightarrow$ builder & $\odot$ & $\odot$ & -- & $\odot$ & -- & $\bullet$ \\
I5\quad builder $\rightarrow$ package index & $\circ$ & $\circ$ & -- & $\odot$ & -- & $\bullet$ \\
I6\quad builder $\rightarrow$ control & $\odot$ & $\bullet$ & -- & $\circ$ & -- & -- \\
I7\quad control $\rightarrow$ executor & $\odot$ & $\bullet$ & -- & $\bullet$ & -- & -- \\
I8\quad executor $\rightarrow$ control & $\odot$ & $\odot$ & -- & $\odot$ & -- & -- \\
\addlinespace
CVM3 $\rightarrow$ CVM1 (reverse pivot) & -- & -- & -- & -- & -- & $\bullet$ \\
CVM3 (per-run hygiene) & -- & -- & -- & -- & $\bullet$ & -- \\
\bottomrule
\end{tabular}
\end{table}

\begin{table}[htbp]
\centering
\caption{Non-trivial STRIDE threats, per interaction of
Figure~\ref{fig:architecture} (S/T/R/I/D/E = spoofing, tampering, repudiation,
information disclosure, denial of service, elevation of privilege). Each row
names the interaction where the threat arises, the countering mechanism, the
measured outcome, and the test providing it; trivial or assumption-covered
cells are discussed in the text.}
\label{tab:stride}
\scriptsize
\begin{tabularx}{\textwidth}{@{}p{1.15cm}cXXp{0.8cm}@{}}
\toprule
Interaction & Cat. & Threat \& countering mechanism & Outcome (measured) & Test \\
\midrule
I1 & S & Caller without access rights; on-chain gate (\texttt{hasAccessByIds}) before any bundle is prepared & Run rejected, \emph{access denied}, before code or data move & T5 \\
\addlinespace
I1, I2 & R & Denial of a registration or purchase; recorded as ledger transactions & Both actions recorded on chain, auditable & T5 \\
\addlinespace
I3 & T & Artifact modified between storage and use; SHA-256 content hash bound on chain & On-chain hash equals file SHA-256 & T5 \\
\addlinespace
I4, I5 & E & Malicious dependency runs at build time; ephemeral builder, filtered egress, no dataset & Payload runs in builder only---no secrets, no dataset; cache write denied & T2, T2b \\
\addlinespace
I6, I7 & T & Poisoned cache or modified artifact; mounted read-only at use & Build-time write denied (\texttt{Permission denied}) & T2b \\
\addlinespace
I7 & I & Dataset exfiltration over the network; execution microVM has no interface & 7/7 outbound channels fail, no \texttt{eth0} & T1, T2a \\
\addlinespace
I7 & I & Reads control-plane secrets or sibling data; no secrets, own disks only & Secrets absent from env; only own I/O disks visible & T1b, T4 \\
\addlinespace
CVM3\,$\to$\,CVM1 & E & Executor pivots into the control plane; secrets and stores absent & Keys, IPFS, cache, executions absent (\emph{residual:} shared net) & T4 \\
\addlinespace
CVM3 & D & Residual state starves later runs; per-run caps (1~vCPU/512~MiB) and teardown & No socket, disk, or mount after 9+ runs & T3 \\
\bottomrule
\end{tabularx}
\end{table}

\textbf{Reading the evidence.} Table~\ref{tab:stride} carries the
per-interaction outcomes; two points deserve emphasis beyond the data. First,
the containment of the execution plane is \emph{structural} rather than
policy-based: the plane that consumes the dataset has no network interface at
all, so the seven outbound probes fail with \texttt{Network is unreachable}
because there is no route to disable, not because a firewall rule denies
them---the strongest form the property can take. This is reinforced by code
size: the execution plane that runs untrusted code is also the smallest
first-party codebase (about 365 lines, against roughly 590 in the builder and
1{,}500 in the secret-holding control plane), so the boundary an attacker
reaches first is the least code to audit and holds none of the control-plane
logic. Second, the supply-chain test (T2) is the paper's thesis in miniature: a
package seeded into the builder's cache is resolved, bundled, and runs its
payload at import time inside the execution microVM---yet finds no secrets and a
blocked network, exactly mirroring T1 and T1b, while its build-time
\texttt{setup.py} variant executes in the isolated builder and fails against the
read-only cache. Malicious code, whether at build or execution time, inherits
the containment of the plane that runs it rather than reaching the control
plane.

\textbf{Trivial, assumption-covered, and residual cells.} The remaining cells
of the sweep in Table~\ref{tab:stride-matrix} fall into three groups.
\emph{Trivial.} On the external-facing interaction I1, tampering and
information disclosure in transit are handled by standard TLS. Interactions
with the ledger and the store (I2, I3) reduce to the trust assumptions of
Section~\ref{sec:proposal}: the ledger is assumed to execute correctly and the
store to be content-addressed, so spoofing and tampering there are outside the
design's control, and the access rights and content hashes they carry are
public by construction. One store cell is a mitigation we do claim rather than
assume: the artifacts are encrypted by the control plane before publication, so
information disclosure at the store (I3) is countered by that encryption.
Spoofing or tampering of a package index (I5) is the classic supply-chain
vector; the design does not authenticate the index beyond standard HTTPS and
instead contains what its packages can do once resolved (the I4/I5 row of
Table~\ref{tab:stride}). Two further trivial cells close the group: information
disclosure of the returned artifact (I6) is immaterial because the control plane
commissioned and already holds it, and elevation of privilege through the
marketplace API itself (I1) is a matter of contract and API implementation
correctness---an assumption outside the architectural analysis rather than a
property the planes enforce. \emph{Acknowledged residuals and out of scope.}
Three residuals are neither trivial nor fully closed in the prototype, and we
mark them as such rather than claim them. First, the internal service channels
(I4, I6--I8) are not mutually authenticated: server-side TLS does not by itself
stop a compromised plane from impersonating a legitimate internal caller, so
spoofing on these channels, in-transit tampering on the build bundle to the
builder (I4), and reach to services that should be unreachable are a
residual that internal mutual TLS and network segmentation would close, the same
network-reach residual already visible in Table~\ref{tab:blast-radius} and
revisited under Threats to Validity. Second, a
malicious dependency running in the builder could exfiltrate the buyer's own
application source over the filtered egress (I4/I5, information disclosure); the
damage is bounded because the provider's dataset never enters the build plane,
so only the buyer's own asset---not the sensitive data---is at risk, but we
record this as an accepted residual rather than a closed threat. Third, the
result returned to the buyer (I8) is out of scope in two senses: this work
bounds what a malicious application can \emph{reach}, not what an authorized
computation \emph{reveals} about its inputs (information disclosure), nor
whether a compromised executor returned a faithful result (tampering);
verifying the processing itself is a complementary line of work discussed in
Section~\ref{sec:related-work}. Denial of service beyond per-run state hygiene
(e.g., request flooding at I1) is likewise out of scope (Threats to Validity),
and repudiation is meaningful only for the ledger-recorded marketplace actions
already covered by the I1/I2 row.

Read against the requirements of Section~\ref{sec:problem-definition}, the
per-interaction analysis exercises R5 (rows I1 and I1/I2 and I3), R3 (rows
I4/I5 and I6), R1 and R4 (both I7 rows and the D row), and R2 (the
CVM3\,$\to$\,CVM1 row); R6 is not a threat-facing property and is exercised
instead by the cross-architecture comparison (Table~\ref{tab:comparison}).

\subsection{Operational Cost}

Table~\ref{tab:comparison} consolidates the cost of the isolation strategy
across architectures, and Table~\ref{tab:ab-performance} reports an earlier
live A/B campaign in detail. End-to-end execution latency is indistinguishable
between A and B; its absolute value tracks host warmth---26.8~s on the
2026-06-11 A/B campaign and 25.2~s on the warmer 2026-06-23 host used here for
the M0 comparison---but the A$\approx$B equality is invariant across both
campaigns. This time is dominated by
the microVM life cycle---kernel boot, disk discovery, output-disk synchronization,
and the guarded shutdown sequence---rather than by the application itself, whose
computation over sixty rows completes in milliseconds.

\begin{table}[htbp]
\centering
\caption{Comparative results across the isolation spectrum: a no-isolation
monolith (M0), a centralized deployment that adds ephemeral microVMs but
colocates all assets (A), the compartmentalized design that also separates the
planes (B), and the SGX-based original DNAT as a cited reference. Execution
latency and overhead are measured on the same host (2026-06-23); on-chain gas
is measured on Hardhat; build and artifact figures, identical across A and B,
are reported with the live A/B campaign (Table~\ref{tab:ab-performance}); the
SGX column is a qualitative reference (its enclave latencies are not
re-measured here).}
\label{tab:comparison}
\footnotesize
\begin{tabularx}{\textwidth}{@{}Xllll@{}}
\toprule
Metric & M0 & A & B & SGX (ref.) \\
\midrule
Execution latency, benign app (s)$^{\P}$ & 0.07 & 25.2 & 25.2 & n/a \\
microVM boot/life-cycle overhead (s) & $\approx$0 & $\approx$25.1 & $\approx$25.1 & n/a \\
On-chain gas, register / purchase (k)$^{\S}$ & 389 / 97 & 389 / 97 & 389 / 97 & --- \\
\addlinespace
App reuse with deps, no recompile & yes & yes & yes & \textbf{no} \\
\bottomrule
\end{tabularx}
\par\smallskip
{\footnotesize The microVM adds $\approx$25~s over M0's $\approx$0.07~s, since the
application's compute over 60 rows is sub-second; that fixed cost is the price of a
fresh, networkless microVM per run. A$\approx$B by design, as both share the same
Firecracker microVMs (build detail in Table~\ref{tab:ab-performance}, blast radius
in Table~\ref{tab:blast-radius}). $^{\P}$The execution figures are the clean-execution comparator (benign app)
from the 2026-06-23 same-host campaign; that campaign's three test apps spanned
25.2--25.8~s. Run-to-run dispersion is characterized in the earlier campaigns
(the $n{=}3$ A/B campaign at 26.8~s in Table~\ref{tab:ab-performance}, and a
colder $n{=}6$ campaign at 31.4~s, sd 1.4~s). On-chain gas is identical
across M0/A/B because the contract is the shared control plane; an asset's optional
bloom filter adds $\approx$185k gas per 256~B. $^{\S}$Gas is the EVM's unit of work
(a plain transfer is 21k); on a representative layer-2 (0.03~gwei, ETH at
US\$3,000) the cost is $\approx$US\$0.035 to register and $<$US\$0.01 to purchase.
The SGX column frames the reuse and dependency limitation of the original DNAT.}
\end{table}

The no-isolation baseline (M0) makes this concrete: running the same workload
directly in the container's namespace---no Firecracker, no disk isolation, no
network removal---the benign application completes in $\approx$0.07~s, so the
microVM adds roughly 25~s of fixed per-execution overhead. That overhead is the
price of launching a fresh, networkless microVM for every run, and it is exactly
what buys the containment: in the same measurement, M0 exposes a working
\texttt{eth0} and succeeds at network exfiltration, finds all six control-plane
secrets in its environment, and persists its writes, whereas A and B (with the
microVM) block exfiltration, expose no secret, and leave no residual state.
Separating the planes (B vs A) adds nothing on top of this microVM cost.

The on-chain cost is likewise architecture-independent. Registering an asset on
the marketplace contract costs $\approx$389k gas (a dataset with an empty bloom
filter; an application $\approx$327k), and purchasing access costs $\approx$97k
gas; a dataset's optional on-chain bloom filter adds $\approx$185k gas per 256
bytes. These figures are identical for M0, A, and B because the registry is the
shared control plane---isolation lives in the build and execution planes, not in
the contract---so the blockchain cost of the user journey is unchanged by the
isolation strategy. In concrete terms (a plain value transfer is 21k gas), on a
representative layer-2 deployment at 0.03~gwei with ETH at US\$3,000 these amount
to about US\$0.035 to register an asset and under US\$0.01 to purchase access, so
the marketplace is inexpensive to operate where such applications realistically
run.

The build plane is more expensive, as expected from its heavier configuration
(two vCPUs, network egress, and dependency resolution), but that cost is
identical across the two designs, which share the same builder microVM. The
per-path build latencies and artifact sizes are reported in detail with the
live A/B campaign in Section~\ref{sec:ab-results}
(Table~\ref{tab:ab-performance}).

\subsection{Threats to Validity}

The measurements were taken on a single host with a Linux-backed container
runtime; absolute latencies will vary with hardware (a colder first campaign saw
a 31.4~s execution mean and 103--131~s builds, with the same A$\approx$B
relationship), but the containment and exfiltration results are structural
(absence of a network interface, absence of secrets) and do not depend on timing.
The execution-plane tests bypass IPFS and the builder to isolate the microVM
under study; the end-to-end path through the control plane is exercised
separately by the access-control test. We report one residual weakness honestly:
in the single-host deployment the three planes share one Docker network, so a
compromised executor can still reach the control API and the blockchain RPC over
that network, even though the secrets and stores that define the blast radius
remain isolated; internal mutual TLS and network segmentation remain future work
(Table~\ref{tab:blast-radius}). Finally, the DoS analysis covers only per-run
state hygiene, not load-induced resource exhaustion, and the SGX column of
Table~\ref{tab:comparison} is qualitative---it restates the original DNAT's
reuse limitation rather than re-measuring enclave latencies here.

A more fundamental caveat concerns the execution substrate. The prototype runs
each plane as a container with Firecracker microVMs rather than inside a
hardware-backed confidential VM, so our measurements do not come from a genuinely
confidential environment and we make no claim about confidential-computing
overheads. This does not weaken the architectural argument: the properties we
evaluate---privilege separation across planes, a networkless execution microVM,
ephemeral teardown, and on-chain access control---are properties of the
architecture, not of the memory-encryption substrate. Running the same planes
inside confidential VMs (e.g., AMD SEV-SNP~\cite{amd2020sevsnp}) would change the
implementation and add a hardware-attestation and memory-encryption cost, but
would preserve the containment and blast-radius results demonstrated here. The
analysis therefore remains valid, with confidential execution as a substrate
substitution rather than an architectural change.

% =====================================================================
% DEPRECATED as an input. The two evaluation tables that used to live here
% (tab:stride and tab:comparison) were moved INTO section_evaluation.tex,
% each right after the paragraph that references it (6.2 STRIDE and 6.3
% Operational Cost). Reason: a LaTeX float never appears earlier than the
% point where it is declared; keeping the tables in a file \input AFTER all
% of section_evaluation's prose forced them to the end of the section.
% The live A/B cost table (tab:ab-performance) and the blast-radius table
% (tab:blast-radius) remain in section_results_ab.tex (subsection 6.5).
% This file is no longer \input from ladc.tex.
% =====================================================================

% =====================================================================
% A/B results: live-measured comparison of Environment A (centralized
% baseline) vs Environment B (compartmentalized DNAT). English, LADC/acmart.
% Self-contained: drop-in tables + a Results subsection that complements
% section_evaluation.tex. Requires booktabs + tabularx (already in preamble).
% Data: TCC/evaluation_data/results_ab_campaign.md (measured 2026-06-11).
% =====================================================================

% ---------------------------------------------------------------------
% Subsection prose: paste at the end of the Evaluation section (or just
% before "Threats to Validity"), after the A column has been promoted
% from "derived from configuration" to "measured live".
% ---------------------------------------------------------------------
\subsection{A/B Comparison: Live Measurement of Both Environments}
\label{sec:ab-results}

To remove any inference from the comparison, we ran the full suite against both deployments
live and end to end: the compartmentalized system (Environment~B) and the centralized
baseline (Environment~A, a single container colocating the same logic). Because both share
the same Firecracker microVMs and the same access gate, the per-microVM guarantees are a
constant---T1, T1b, T3 and T5 reproduce in both, and so does the operational cost---so the
architectural variable surfaces only in the blast radius. That is where A and B diverge
(Table~\ref{tab:blast-radius}, populated from live inspection rather than configuration):
in B, the executor and the builder carry \emph{none} of the control-plane secrets, and the
sensitive stores (IPFS repository, wheel cache, executions store) are absent from their
namespace, whereas in A those same secrets and stores sit in the single namespace that also
runs the execution logic. The supply-chain test (T2) inherits the split: a compromised
dependency's code executes, but in B inside a plane that holds neither secrets nor a route
to them.

\textbf{What the network residual still grants.} Reachability is not exploitability, so we
state what a compromised executor can \emph{do} with the two ports it still reaches in B.
The prototype does not authenticate its internal endpoints, so the control API can be driven
as a confused deputy: an attacker on the executor can enumerate registered assets, read the
stored output of previous executions, and ask the control plane to act with its wallet
through the fixed set of operations the API exposes. This is a real residual, and it is why
internal mutual TLS heads our future work. What it does \emph{not} give is the platform: the
encryption and wallet keys are never returned by any endpoint and remain in CVM1's
environment, so the attacker can neither decrypt artifacts pulled from the store nor sign
transactions of its own choosing, and the on-chain gate still refuses an execution for which
no access was purchased (T5). In A the same attacker needs no network hop at all---both keys
sit in its own environment and the stores on its filesystem---so it can decrypt every
archived artifact and sign arbitrary transactions directly. The difference is therefore not
the presence of a residual but its ceiling: a bounded, mediated set of API calls in B
against unmediated custody of the platform's keys in A---\emph{containment, not absolute
security}.

% ---------------------------------------------------------------------
% Table: blast radius B vs A — English version, live-measured wording.
% A/B cost figures live in tab:comparison (section_evaluation.tex).
% ---------------------------------------------------------------------
\begin{table}[htbp]
\centering
\caption{Reach of a compromised execution plane over control-plane assets: the
compartmentalized design (B) versus the centralized baseline (A), from live inspection.
``Isolated'' = absent from the plane's namespace (the environment variable is not set, the
store is not mounted); ``Colocated'' = same namespace as the execution logic.}
\label{tab:blast-radius}
\footnotesize
\begin{tabularx}{\textwidth}{@{}>{\raggedright\arraybackslash}X>{\raggedright\arraybackslash}p{2.4cm}>{\raggedright\arraybackslash}p{2.4cm}@{}}
\toprule
Control-plane asset reachable from the execution plane & Env.~B (proposed) & Env.~A (baseline) \\
\midrule
Asset-encryption key (env) & Isolated & Colocated \\
Wallet private key (env) & Isolated & Colocated \\
Service URLs: RPC, IPFS, builder, executor (env) & Isolated & Colocated \\
IPFS repository, wheel cache, executions store (filesystem) & Isolated & Colocated \\
IPFS API (port 5001, network) & Blocked & Reachable \\
Control API (3001) and blockchain RPC (8545), network & Reachable$^{\dagger}$ & Reachable \\
\bottomrule
\end{tabularx}
\par\smallskip
\parbox{\textwidth}{\footnotesize $^{\dagger}$Residual risk in the single-host deployment:
the three containers share one Docker network and the internal endpoints are
unauthenticated, so a compromised executor can invoke the control API as a confused deputy
(enumerate assets, read stored execution outputs) without ever holding its keys---see the
discussion above. The defining assets of the blast radius---secrets and persistent
stores---are nonetheless isolated in B and colocated in A.}
\end{table}

% ---------------------------------------------------------------------
% Threats to Validity: movida de section_evaluation.tex (texto intacto)
% para fechar a secao de Evaluation depois da subsecao A/B.
% ---------------------------------------------------------------------
\subsection{Threats to Validity}

The measurements were taken on a single host; absolute latencies will vary with hardware (a
colder first campaign saw a 31.4~s execution mean and 103--131~s builds, with A and B again
indistinguishable), and each cost figure is a single run, so we read them as orders of
magnitude rather than precise separations. The containment and exfiltration results, by
contrast, are structural---absence of a network interface, absence of secrets---and do not
depend on timing. Our workload is also deliberately small: a 60-row dataset isolates the
fixed microVM cost, but says nothing about how that cost amortizes over a long-running job,
nor about where the per-run caps (1~vCPU, 512~MiB, fixed disk sizes) begin to bind. The
network-reach residual bounds the blast-radius claim, with the consequences and limits set
out above, and confidentiality against the infrastructure operator (the second half of R1)
is assumed rather than shown, for the substrate reason given next. Finally, the DoS analysis
covers only per-run state hygiene, and the SGX column of Table~\ref{tab:comparison} is
qualitative rather than re-measured.

A more fundamental caveat concerns the execution substrate. The prototype runs each plane as
a container with Firecracker microVMs rather than inside a hardware-backed confidential VM,
so our measurements do not come from a genuinely confidential environment and we make no
claim about confidential-computing overheads. This does not weaken the architectural
argument: the properties we evaluate---privilege separation across planes, a networkless
execution microVM, ephemeral teardown, and on-chain access control---belong to the
architecture, not to the memory-encryption substrate. Running the same planes inside
confidential VMs (e.g., AMD SEV-SNP~\cite{amd2020sevsnp}) would add attestation and
memory-encryption costs---so the cost figures above would not carry over unchanged---but
would preserve the containment and blast-radius results, which is why we treat confidential
execution as a substrate substitution rather than an architectural change.

\clearpage


\section{Related Work}\label{sec:related-work}

% mencionar o DNAT e suas limitações (ex., dependência de tecnologias limitadas)

% ver os artigos que ele citou ou que citaram ele

\section{Concluding Remarks}\label{sec:concluding-remarks}

% =====================================================================
% Conclusao revisada, integrando as evidencias experimentais.
% Substitui a Secao "Conclusion" atual do main.tex. Idioma: ingles.
% =====================================================================

\section{Conclusion}
\label{sec:conclusion}

This paper framed the proposed system as a layered containment architecture for a
confidential data marketplace and tested that framing against the running prototype. The
central design move---separating control, build, and execution into planes with different
network and persistence properties---was evaluated not as an abstract claim of trustlessness
but as a measurable reduction in blast radius relative to a centralized baseline differing
only in the colocation of assets.

The experimental results support that argument. Execution-time code cannot open outbound
channels: the probed channels failed in the network stack itself---unimplemented socket
calls and failed name resolution---rather than at a policy that could be switched off.
Ephemeral state does not survive, and the assets that define the blast radius (the
encryption and wallet keys, the IPFS repository, the wheel cache, and the executions store)
are absent from the compromised executor in the compartmentalized design, whereas they are
colocated in the baseline. Execution is gated on chain, and the isolation costs roughly
twenty-five seconds per run---a cost compartmentalization itself does not increase.

The limitations are equally concrete: the planes share one network and do not authenticate
each other, so a compromised executor can still drive the control API as a confused
deputy---though never obtaining the keys to decrypt artifacts or sign transactions; the
wheel cache is not yet a trustworthy artifact store; dataset confidentiality depends on
preparation outside the encrypted-artifact path; and no attestation flow yet verifies each
component's software identity. The design is nonetheless a credible high-assurance baseline,
and future work---internal mutual TLS, wheel and dataset provenance, and attestation---will
let the containment shown here be verified remotely rather than assumed.


% Repetir alguns pontos chave da motivação e como a arquitetura resolve. E colocar algo na linha de número de desempenho para mostrar que é viável.

% Mencionar possíveis trabalhos futuros, ex.: que blockchain (custo, escalabilidade, atrasos)? IPFS? Outras linguagens? etc.

\bibliographystyle{splncs04}
\bibliography{sample}

\end{document}
